%! AP Statistics — Unit 8, Lesson 8.2 — Homework (Answer Key)
\documentclass[10pt]{article}
\usepackage{apstats-article}
\usepackage{apstats-key}
\usepackage{grading}
\gradingassignment{apstats-u08-l02-hw}
\gradingcourse{advanced_statistics}

\begin{document}

\pageheader{Unit 8, Lesson 8.2}{Homework --- Answer Key}
\namedateperiod

\vspace{0.1in}
\textbf{Directions:} For each problem, set up a chi-square goodness of fit test as directed.
Show all work clearly.

\begin{enumerate}

  \item \textbf{Genetics.} $n = 160$ offspring; 9:3:3:1 ratio.

        \begin{enumerate}[label=\alph*.]
          \item H$_0$: \ansline{$p_A = 9/16$, $p_B = 3/16$, $p_C = 3/16$, $p_D = 1/16$
                (phenotypes follow the 9:3:3:1 Mendelian ratio).}\gradedkeyline{\gradedblank[id=hw1a1, type=rubric, points=1, display={$p_A = 9/16$, $p_B = p_C = 3/16$, $p_D = 1/16$ (9:3:3:1 ratio holds)}]{\linewidth-1.0in}{\gradedlinesheight{2}}}

                H$_a$: \ansline{At least one phenotype proportion is not as specified by the 9:3:3:1 ratio.}\gradedkeyline{\gradedblank[id=hw1a2, type=rubric, points=1, display={At least one proportion differs from 9:3:3:1.}]{\linewidth-1.0in}{\gradedlinesheight{1}}}

          \item \begin{center}
                \renewcommand{\arraystretch}{1.4}
                \begin{tabular}{lcccc}
                    \textbf{Phenotype} & A & B & C & D \\
                    \hline
                    Null proportion & $9/16$ & $3/16$ & $3/16$ & $1/16$ \\
                    \hline
                    Expected count & \makebox[2cm][r]{\smash{\raisebox{-3.5pt}{\gradedblank[id=hw1b1, type=short, points=1, answer=90, check=numeric, tolerance=0]{0.8cm}{0.2in}}}} & \makebox[2cm][r]{\smash{\raisebox{-3.5pt}{\gradedblank[id=hw1b2, type=short, points=1, answer=30, check=numeric, tolerance=0]{0.8cm}{0.2in}}}} & \makebox[2cm][r]{\smash{\raisebox{-3.5pt}{\gradedblank[id=hw1b3, type=short, points=1, answer=30, check=numeric, tolerance=0]{0.8cm}{0.2in}}}} & \makebox[2cm][r]{\smash{\raisebox{-3.5pt}{\gradedblank[id=hw1b4, type=short, points=1, answer=10, check=numeric, tolerance=0]{0.8cm}{0.2in}}}} \\
                \end{tabular}
                \end{center}

          \item \ansline{All expected counts $\ge 5$ (smallest is 10). Large Counts condition is met.}\gradedkey{\gradedblank[id=hw1c, type=rubric, points=2, display={Smallest expected 10 $\ge 5$: met.}]{2.6in}{0.19in}}
        \end{enumerate}

  \item \textbf{Transportation Survey.} $n = 200$; four modes, each claimed at 25\%.

        \begin{enumerate}[label=\alph*.]
          \item H$_0$: \ansline{$p_\text{car} = p_\text{bus} = p_\text{bike} = p_\text{walk} = 0.25$
                (commuters use the four modes in equal proportions).}\gradedkeyline{\gradedblank[id=hw2a1, type=rubric, points=1, display={$p_\text{car} = p_\text{bus} = p_\text{bike} = p_\text{walk} = 0.25$}]{\linewidth-1.0in}{\gradedlinesheight{1}}}

                H$_a$: \ansline{At least one transportation mode is used by a proportion of commuters
                that is not 0.25.}\gradedkeyline{\gradedblank[id=hw2a2, type=rubric, points=1, display={At least one mode's proportion is not 0.25.}]{\linewidth-1.0in}{\gradedlinesheight{1}}}

          \item $E = 200 \times 0.25 = $ \ans{50 for each mode.}\gradedkey{\gradedblank[id=hw2b, type=short, points=1, answer=50, check=numeric, tolerance=0]{6cm}{0.17in}}

                \begin{center}
                \renewcommand{\arraystretch}{1.3}
                \begin{tabular}{lcccc}
                    Mode & Car & Bus & Bike & Walk \\
                    \hline
                    Expected & \ans{50} & \ans{50} & \ans{50} & \ans{50} \\
                \end{tabular}
                \end{center}

          \item \begin{itemize}
                    \item \textbf{Random:} \ans{Random sample of 200 commuters stated. \checkmark}\gradedkey{\gradedblank[id=hw2c1, type=rubric, points=1, display={Random sample stated.}]{7cm}{0.19in}}
                    \item \textbf{Large Counts:} \ans{All expected counts $= 50 \ge 5$. \checkmark}\gradedkey{\gradedblank[id=hw2c2, type=rubric, points=1, display={All expected $= 50 \ge 5$.}]{7cm}{0.19in}}
                \end{itemize}
        \end{enumerate}

  \item \textbf{True or False.}

        \textcolor{keyred}{\textbf{False.}}\gradedkey{\textit{T/F:}~\gradedblank[id=hw3a, type=choice, points=1, answer=F]{0.4in}{0.2in}}

        \ansline{We check that each \emph{expected} count (not observed count) is $\ge 5$.
        The expected counts are calculated from the null hypothesis: $E = n \cdot p_0$.
        Even with a large sample, an observed count being large does not guarantee the
        expected count is $\ge 5$ --- and the condition applies to expected, not observed, counts.}\gradedkeyline{\gradedblank[id=hw3b, type=rubric, points=2, display={False: the condition is on \emph{expected} counts ($n p_0 \ge 5$), not observed counts.}]{\linewidth+6pt}{\gradedlinesheight{2}}}

\end{enumerate}

\end{document}
